Lo que más llama la atención es lo rápido que se enfocan los rayos. Con una sola superficie curva de acrílico el foco estaría a más de tres radios, pero en el cilindro los rayos se cruzaron prácticamente pegados a la pieza, y la razón es que cruzan dos superficies curvas y no una: la de entrada empieza a juntarlos, la de salida los vuelve a doblar en la misma dirección y, sumando el efecto de las dos, el foco queda a $1.52\,R$ del centro, como indica la ecuación \eqref{eq:cilindro}.

En otras palabras, el cilindro se comporta como una lente convergente muy potente.

El segundo resultado es la aberración esférica. Que los rayos de las orillas se crucen antes que los del centro es exactamente lo que se espera, porque llegan a la superficie con ángulos de incidencia grandes, donde la aproximación paraxial ya no funciona y la ecuación \eqref{eq:snell} los dobla más de lo que predice la ecuación \eqref{eq:superficie}. En el cilindro este efecto es muy notorio porque los rayos de los extremos entran lejos del eje en comparación con el radio de la pieza. Por eso las lentes de calidad no usan superficies tan curvas, o combinan varias lentes para corregirlo, y por eso también las cámaras dan imágenes más nítidas cuando se cierra el diafragma y solo pasan los rayos centrales.

Los colores tienen otra explicación. El acrílico no dobla igual todos los colores, y como la fuente da luz blanca, cada componente se enfoca en un lugar un poquito distinto. Este defecto se llama aberración cromática.

Hay que reconocer que, al no haber mediciones, la comparación con la teoría es solo cualitativa. Habría sido útil medir el radio de la pieza y la distancia del cruce de los rayos centrales hasta la cara de salida para comprobar el valor de $0.52\,R$, así como la distancia de cruce de los rayos de las orillas para tener una idea de qué tan grande es la aberración.

En conclusión, una superficie curva de acrílico hace converger un haz de rayos paralelos, y un cilindro completo, al tener dos superficies, lo enfoca a muy poca distancia de su cara de salida, en acuerdo con la ecuación de una superficie refractante esférica aplicada dos veces. Además se observaron con claridad dos limitaciones de las lentes reales: la aberración esférica, por la que los rayos de las orillas se enfocan antes que los centrales, y la aberración cromática, por la que cada color se enfoca en un punto diferente.
